\documentclass[11pt]{article}

% ====================================================
% PACKAGES
% ====================================================
\usepackage[margin=1in, headheight=13.6pt]{geometry}
\usepackage[T1]{fontenc}
\usepackage{lmodern}
\usepackage{microtype}
\usepackage{setspace}
\usepackage{amsmath, amssymb, mathtools}
\usepackage{siunitx}

\usepackage{graphicx}
\usepackage{float}
\usepackage{subcaption}

\usepackage{enumitem}
\usepackage{booktabs}
\usepackage{longtable}
\usepackage{array}

\usepackage{xcolor}
\usepackage{hyperref}

\usepackage{fancyhdr}
\usepackage{lastpage}

\usepackage{indentfirst}
\usepackage[justification=centering]{caption}

\pagestyle{fancy}
\fancyhf{}
\lhead{\Assignment\ |\ \course}
\rhead{\Name\ |\ \today}
\cfoot{Page \thepage\ of~\pageref{LastPage}}

\hypersetup{
    colorlinks,
    linkcolor={red!50!black},
    citecolor={blue!50!black},
    urlcolor={blue!80!black}
}

\newcommand{\Name}{Arael Anaya}
\newcommand{\Assignment}{3.5 Link Budget Analysis}
\newcommand{\course}{Space Robotics}

\begin{document}

% ====================================================
% SECTION 0: LINK DEFINITION
% ====================================================
\section*{I. Link Definition and Desired Margin}

\begin{enumerate}[label=\alph*., leftmargin=*]
    \item \textbf{Next communication node:} the nearer of the two orbital relay agents. The rover's
    mission doesn't route through a direct link to Earth; Section~VI works out why. This report
    follows the assignment's naming: rover~$\rightarrow$~relay counts as uplink, since that's the
    direction the science data comes back on, and relay~$\rightarrow$~rover counts as downlink,
    carrying commands out to the rover.

    \item \textbf{Band and channel:} UHF proximity link on the CCSDS Proximity-1 return/forward
    pair, \SI{401.6}{MHz} return and \SI{437.1}{MHz} forward, the same split flown by Electra and
    Electra-Lite on the Mars Relay Network.

    \item \textbf{Desired margin:} \SI{3}{dB} above the $E_b/N_0$ required for a post-decoding bit
    error rate of $10^{-6}$. CCSDS turbo coding at rate $1/2$ puts that threshold near
    \SI{1.5}{dB}, so the link has to clear \SI{4.5}{dB} total. Three decibels is the right cushion
    given how much is still uncertain before flight. Nobody has flown the relay's exact orbit
    geometry yet. The gain pattern on the rover's helix is still just a model; nobody has
    range-tested the flight unit. And Jupiter's synchrotron emission adds an amount of receiver
    noise that is still a rough guess at this stage.

    \item \textbf{Worst-case geometry:} both relays fly a \SI{200}{km} circular Europa orbit. With
    $R_{\mathrm{Europa}} = \SI{1561}{km}$ and a \SI{10}{\degree} elevation mask, the slant range at
    the edge of the pass works out to
    \[
        R = -R_E \sin\varepsilon + \sqrt{R_E^2\sin^2\varepsilon + \left(r^2 - R_E^2\right)}
          = \SI{588}{km},
    \]
    compared with \SI{200}{km} at zenith. Every budget from here on uses the \SI{588}{km} worst case.
\end{enumerate}

% ====================================================
% SECTION II: ASSUMED PARAMETERS
% ====================================================
\section*{II. Assumed Parameters}

\begin{enumerate}[label=\alph*., leftmargin=*]
    \item $P_t = \SI{5}{W}$ RF (rover), $\SI{8}{W}$ RF (relay). The rover's UHF radio draws
    \SI{35}{W} DC in the power budget, and an Electra-Lite class transceiver puts about
    \SI{5}{W} of that onto the antenna. The relay's power budget has more slack than the rover's,
    so it carries a larger amplifier.

    \item $G_t = 3.16$ (\SI{5.0}{dBi}) rover quadrifilar helix; $G_r = 5.01$ (\SI{7.0}{dBi}) relay
    UHF patch array. The helix already has a spot in the mass budget, and its near-hemispherical
    pattern means the rover can close a pass without pointing at anything.

    \item $L_a = 0.89$ ($\SI{-0.51}{dB}$). Europa's atmosphere, at roughly $10^{-12}$~bar, is far
    too thin to absorb anything here. This loss term comes from circular-to-linear polarization
    mismatch and Jovian plasma scintillation along the path.

    \item $L_l = 0.71$ ($\SI{-1.49}{dB}$), the transmitter-to-antenna line loss. It comes from the
    cabling and connectors running out to a mast-mounted antenna, derated for operation at a
    cryogenic \SI{90}{K}.

    \item $\lambda = \SI{0.7465}{m}$ return (\SI{401.6}{MHz}), $\lambda = \SI{0.6859}{m}$ forward
    (\SI{437.1}{MHz}), from $\lambda = c/f$.

    \item $k = \SI{1.38e-23}{J/K}$.

    \item $T_s = \SI{3000}{K}$ at the relay receiver, $\SI{4000}{K}$ at the rover receiver. This is
    the number that really separates a Europa proximity link from a Mars one: at \SI{400}{MHz} the
    receiver picks up galactic background noise on top of Jupiter's synchrotron emission, and the
    rover's front end, with no cryogenic cooling and a warmer view of the sky than the relay's,
    ends up carrying the higher figure.

    \item $R = \SI{5.879e5}{m}$, the \SI{10}{\degree}-elevation slant range derived above.

    \item $R_b = \SI{128}{kbps}$ return, \SI{8}{kbps} forward. REQ.10.01 sets the floor here:
    getting \SI{55}{MB} of compressed data down inside the \SI{90}{min} relay window the ConOps
    allocates takes about \SI{90}{kbps}, and \SI{128}{kbps} is the next standard rate above that.
\end{enumerate}

% ====================================================
% SECTION III: METHOD
% ====================================================
\section*{III. Method}

\begin{enumerate}[label=\alph*., leftmargin=*]
    \item Ratio form:
    \[
        \frac{E_b}{N_0} = \frac{P_t G_t G_r L_a L_l \lambda^2}{(4\pi R)^2\, k\, T_s\, R_b}
    \]
    \item Decibel form: $\left(E_b/N_0\right)_{\mathrm{dB}} = 10\log_{10}\left(E_b/N_0\right)$
    \item Margin: $M = \left(E_b/N_0\right)_{\mathrm{dB}} - \left(E_b/N_0\right)_{\mathrm{req}}$,
    with $\left(E_b/N_0\right)_{\mathrm{req}} = \SI{1.5}{dB}$.
    \item The link closes when $M \geq \SI{3}{dB}$.
\end{enumerate}

% ====================================================
% SECTION IV: UPLINK
% ====================================================
\section*{IV. Uplink: Rover to Relay (Science Return)}

\begin{enumerate}[label=\alph*., leftmargin=*]
    \item Numerator:
    \[
        P_t G_t G_r L_a L_l \lambda^2
        = 5 \times 3.16 \times 5.01 \times 0.89 \times 0.71 \times (0.7465)^2
        = 27.87
    \]
    \item Denominator:
    \[
        (4\pi R)^2 k T_s R_b
        = \left(4\pi \times 5.879\!\times\!10^5\right)^2 \times \SI{1.38e-23}{} \times 3000
          \times 1.28\!\times\!10^5
        = 0.2892
    \]
    \item Result: $E_b/N_0 = 27.87 / 0.2892 = 96.4$
    \item In decibels: $10\log_{10}(96.4) = \SI{19.84}{dB}$
    \item Margin: $19.84 - 1.5 = \SI{18.3}{dB}$. \textbf{The uplink closes.}
\end{enumerate}

\noindent Here's the same result broken down term by term, in decibels, so each contribution is
traceable:

\begin{center}
\begin{tabular}{lr}
\toprule
Term & Value (dB) \\
\midrule
$P_t$ (\SI{5}{W})                        & $+6.99$ \\
$G_t$ (rover helix)                      & $+5.00$ \\
$G_r$ (relay patch array)                & $+7.00$ \\
$L_a$ (polarization + plasma)            & $-0.51$ \\
$L_l$ (line loss)                        & $-1.49$ \\
$\lambda^2$                              & $-2.54$ \\
$(4\pi R)^2$ (space loss at \SI{588}{km})& $-137.37$ \\
$k$                                      & $+228.60$ \\
$T_s$ (\SI{3000}{K})                     & $-34.77$ \\
$R_b$ (\SI{128}{kbps})                   & $-51.07$ \\
\midrule
$E_b/N_0$                                & $\mathbf{+19.84}$ \\
\bottomrule
\end{tabular}
\end{center}

% ====================================================
% SECTION V: DOWNLINK
% ====================================================
\section*{V. Downlink: Relay to Rover (Command Forward)}

\begin{enumerate}[label=\alph*., leftmargin=*]
    \item Numerator:
    \[
        8 \times 5.01 \times 3.16 \times 0.89 \times 0.71 \times (0.6859)^2 = 37.65
    \]
    \item Denominator:
    \[
        \left(4\pi \times 5.879\!\times\!10^5\right)^2 \times \SI{1.38e-23}{} \times 4000
          \times 8.0\!\times\!10^3 = 0.02410
    \]
    \item Result: $E_b/N_0 = 37.65 / 0.02410 = 1562$
    \item In decibels: $10\log_{10}(1562) = \SI{31.94}{dB}$
    \item Margin: $31.94 - 1.5 = \SI{30.4}{dB}$. \textbf{The downlink closes.}
\end{enumerate}

\noindent The forward direction has it easier. The relay pushes more power out of a higher-gain
antenna, and it only has to move \SI{8}{kbps} of commands against \SI{128}{kbps} of science data
coming back the other way. That \SI{16}{\times} difference in rate by itself is worth \SI{12}{dB}
of the gap.

% ====================================================
% SECTION VI: DIRECT-TO-EARTH CHECK
% ====================================================
\section*{VI. Direct-to-Earth Check (Why the Relay Exists)}

Running the same equation on the rover's X-band low-gain antenna against a DSN \SI{34}{m} BWG
station gives the fallback rate. With $P_t = \SI{15}{W}$, $G_t = 6.31$ (\SI{8}{dBi} LGA),
$G_r = 6.6\!\times\!10^6$ (\SI{68.2}{dBi}), $\lambda = \SI{0.0357}{m}$ (\SI{8.4}{GHz}), and
$T_s = \SI{30}{K}$ (a cryogenically cooled DSN front end), solving for the rate that still holds
\SI{4.5}{dB} gives:

\begin{center}
\begin{tabular}{lrr}
\toprule
Earth range & $R$ (m) & Max rate at \SI{4.5}{dB} \\
\midrule
Opposition ($\SI{4.2}{AU}$)  & $6.28\!\times\!10^{11}$ & \SI{6.9}{bps} \\
Conjunction ($\SI{6.2}{AU}$) & $9.27\!\times\!10^{11}$ & \SI{3.2}{bps} \\
\bottomrule
\end{tabular}
\end{center}

\noindent Seven bits per second covers a state-of-health beacon. One \SI{55}{MB} sol of REQ.10.01
data would take just over two years to clear at that rate, which is why the relay agents belong in
the design from the start. REQ.10.03 fixes this same number as the \SI{5}{bps} fallback floor.

% ====================================================
% SECTION VII: REFLECTION
% ====================================================
\section*{VII. Reflection}

Both directions works, and by a wide margin: \SI{18.3}{dB} on the return, \SI{30.4}{dB} on the
forward. The proximity link seems to run out of headroom somewhere other than the RF budget.
At \SI{4.5}{dB} the return could carry something like \SI{4.4}{Mbps}, well past what the hardware
and the pass geometry allow: the radio tops out at \SI{2048}{kbps}, and the pass itself only lasts
\SI{90}{min}. That surplus is worth utilizing, since $T_s$ is the weakest assumption in
this analysis, and Jupiter's synchrotron noise at \SI{400}{MHz} could run several times
higher than the \SI{3000}{K} figure I used. If the link failed anyway, $R_b$ is the first thing I'd
change, since adjusting it costs nothing. $G_r$ at the relay comes next. Adding patch elements to
an orbiter costs less in mass and power than any equivalent change on the rover.

% ====================================================
% REFERENCES
% ====================================================
\bigskip

\noindent\underline{\textbf{References:}}

\begin{enumerate}[leftmargin=*, label={[\arabic*]}]
    \item Jet Propulsion Laboratory, ``DSN Telecommunications Link Design Handbook,'' 810-005,
    Rev.~E, NASA/JPL, Pasadena, CA. [Online]. Available:
    \url{https://deepspace.jpl.nasa.gov/dsndocs/810-005/}.
    \item Consultative Committee for Space Data Systems, ``Proximity-1 Space Link Protocol ---
    Physical Layer,'' CCSDS 211.1-B-4, 2013.
    \item Consultative Committee for Space Data Systems, ``TM Synchronization and Channel Coding,''
    CCSDS 131.0-B-4, 2017.
    \item C.~D.~Edwards \textit{et al.}, ``The Electra Proximity Link Payload for Mars Relay
    Telecommunications and Navigation,'' Jet Propulsion Laboratory, Pasadena, CA.
    \item NASA, ``Deep Space Network,'' NASA Space Communications and Navigation. [Online].
    Available: \url{https://www.nasa.gov/communicating-with-missions/dsn/}.
    \item NASA, ``Radar for Europa Assessment and Sounding: Ocean to Near-surface (REASON),''
    Europa Clipper Mission. [Online]. Available:
    \url{https://europa.nasa.gov/spacecraft/instruments/reason/}.
\end{enumerate}

\bigskip

\noindent{\color{red}\textit{AI disclosure: I used AI to check the arithmetic behind each link
budget term and to cross-check the assumed parameters against published Electra and DSN
performance data. I picked the architecture and the parameters myself, and I set the margin the
same way.}}

\end{document}